To isolate reference-model bias, we replace Pythia-410M with two alternative references from different families and evaluate on a 5-generator subset (Pythia-2.8B/1B, OLMo-2-1B, Qwen3-0.6B/1.7B).
Table~\ref{tab:cross-ref} and Figure~\ref{fig:ref-comparison} show that performance spread across references is only 0.013 F1, confirming the DDR partition is reference-model agnostic.

\begin{table}[t]
\centering
\footnotesize
\renewcommand{\arraystretch}{1.3}
\caption{Cross-family reference model validation (5-gen LOGO subset). Drift direction is consistent (negative) across all three references; magnitude varies with capacity. Spread in LOGO F1 is 0.013.}
\label{tab:cross-ref}
\resizebox{\linewidth}{!}{%
\begin{tabular}{llccc}
\toprule
\textbf{Reference} & \textbf{Family} & \textbf{Params} & \textbf{LOGO F1} & \textbf{Ent.\ Drift} \\
\midrule
Pythia-410M & Pythia & 410M & 0.803 & [$-$1.19, $-$1.40] \\
OLMo-2-1B  & OLMo   & 1.5B & 0.807 & [$-$0.61, $-$0.91] \\
Qwen3-0.6B & Qwen   & 0.6B & 0.816 & [$-$0.96, $-$1.35] \\
\bottomrule
\end{tabular}}
\end{table}

\begin{figure}[t]
\centering
\includegraphics[width=0.85\columnwidth]{fig_appendix_ref_comparison.pdf}
\caption{Cross-family reference model comparison. LOGO F1 spread is only 0.013 across three reference models from different families, confirming DDR is reference-model agnostic.}
\label{fig:ref-comparison}
\end{figure}
